\documentclass[level=3]{kidscode}
\usepackage{multicol}
%% Compact headings, code and spacing so the sheet fits on two sides.
\titleformat{\subsection}{\needspace{5\baselineskip}\normalsize\bfseries\color{kcaccent}}{}{0pt}{}
\titlespacing*{\subsection}{0pt}{5pt}{1pt}
\lstdefinestyle{kcpython}{
  language=Python,
  basicstyle=\ttfamily\footnotesize\setstretch{1.0},
  keywordstyle=\bfseries\color{kcblue},
  stringstyle=\color{kcvermillion!85!black},
  commentstyle=\color{kcgrey},
  showstringspaces=false, columns=fullflexible, keepspaces=true,
  breaklines=true, numbers=none, tabsize=4, upquote=true,
  morekeywords={True,False,None,self}}
%% Plain framed listings break cleanly inside multicols (tcolorbox does not).
\lstnewenvironment{sheetcode}{\lstset{style=kcpython,frame=single,framesep=3pt,
  rulecolor=\color{kcsky!70!black},backgroundcolor=\color{kcsky!7},
  xleftmargin=3.5pt,xrightmargin=3.5pt,aboveskip=3pt,belowskip=4pt}}{}
\setlength{\columnsep}{16pt}
\setlength{\parskip}{0.3em}
\raggedcolumns

\begin{document}
\kctitle{Cheat Sheet}{Classes, algorithms, data and git on two pages}
\small
\begin{multicols}{2}

\subsection{Classes and objects}
\begin{sheetcode}
class Pet:                # blueprint
    def __init__(self, name):
        self.name = name  # attribute
        self.energy = 5

    def play(self):       # method
        self.energy -= 2
        return self.name + " had fun!"

mo = Pet("Mo")            # an object
print(mo.play(), mo.energy)
\end{sheetcode}
\textbf{Attributes} = what it has. \textbf{Methods} = what it does. \texttt{self} means
\enquote{this object}; every method needs it first.

\subsection{Inheritance (IS a), composition (HAS a)}
\begin{sheetcode}
class Dragon(Pet):        # IS a Pet
    def __init__(self, name):
        super().__init__(name)
        self.fire = 3

    def play(self):       # override
        return "Roar!"

class Zoo:                # HAS pets
    def __init__(self):
        self.pets = [Pet("Mo"),
                     Dragon("Ash")]

    def play_all(self):
        for pet in self.pets:
            print(pet.play())
\end{sheetcode}

\subsection{Searching: how fast is fast?}
\begin{sheetcode}
def binary_search(items, target):
    # items must be SORTED
    low, high = 0, len(items) - 1
    while low <= high:
        mid = (low + high) // 2
        if items[mid] == target:
            return mid    # found it
        elif items[mid] < target:
            low = mid + 1  # right
        else:
            high = mid - 1 # left
    return -1             # not there
\end{sheetcode}
\begin{tabular}{@{}lrrr@{}}
\toprule
Items & 10 & 1000 & 1\,000\,000\\ \midrule
Linear search, at most & 10 & 1000 & 1\,000\,000\\
Binary search, at most & 4 & 10 & 20\\
\bottomrule
\end{tabular}

\subsection{Sorting}
\textbf{Selection sort:} find the smallest, swap it to the front, repeat with the rest.
\textbf{Bubble sort:} swap neighbours in the wrong order; repeat until nothing swaps.
\begin{sheetcode}
print(sorted([7, 3, 9]))  # [3, 7, 9]

def get_score(player):
    return player[1]

players = [("Mei", 40), ("Leo", 75)]
best = sorted(players, key=get_score,
              reverse=True)  # Leo 1st
\end{sheetcode}

\subsection{Stack (last in, first out)}
\begin{sheetcode}
stack = []            # undo, plates
stack.append("circle")
stack.append("square")
last = stack.pop()    # "square"
\end{sheetcode}

\subsection{Queue (first in, first out)}
\begin{sheetcode}
from collections import deque
queue = deque()       # turns, BFS
queue.append("Ada")
queue.append("Kai")
first = queue.popleft()   # "Ada"
\end{sheetcode}

\subsection{Recursion}
A function that calls itself on a \textbf{smaller} problem. It needs a \textbf{base case} to stop,
or Python gives a \texttt{RecursionError}.
\begin{sheetcode}
def tree(t, length):
    if length < 10:       # base case
        return
    t.forward(length)
    t.left(30)
    tree(t, length * 0.7) # smaller
    t.right(60)
    tree(t, length * 0.7)
    t.left(30)
    t.backward(length)    # go back
\end{sheetcode}

\subsection{Grids (2D lists)}
\begin{sheetcode}
grid = [["."] * 5 for r in range(3)]
grid[1][2] = "#"          # [row][col]
for row in grid:
    print("".join(row))
\end{sheetcode}
Neighbours of \texttt{(r, c)}: up \texttt{(r-1, c)}, down \texttt{(r+1, c)}, left
\texttt{(r, c-1)}, right \texttt{(r, c+1)}.

\textbf{BFS (shortest path):} put the start in a \texttt{deque}. Repeat: \texttt{popleft()} a square, add
its unvisited neighbours to the back, remember where each came from.

\subsection{CSV data}
\begin{sheetcode}
import csv
with open("games.csv", newline="",
          encoding="utf-8") as f:
    for row in csv.DictReader(f):
        text = row["rating"].strip()
        if text == "":
            continue      # skip blanks
        rating = int(text)  # text->int

# count with a dict:
counts[x] = counts.get(x, 0) + 1
\end{sheetcode}

\subsection{git: save points (in Terminal)}
\renewcommand{\arraystretch}{1.1}
\begin{tabular}{@{}>{\ttfamily}l>{\raggedright\arraybackslash}p{3.3cm}@{}}
git init & start a repo (once)\\
git status & what changed?\\
git diff & show changed lines\\
git add game.py & choose a file (\texttt{.} = all)\\
git commit -m "Add X" & make a save point\\
git log --oneline & list save points\\
git restore game.py & undo changes not yet committed\\
pwd & where am I?\\
cd folder & go into a folder\\
\end{tabular}

\textbf{Good message:} verb + what changed, e.g.\ \texttt{"Fix ghost walking through walls"}.
Never commit secrets or real names.

\subsection{pytest}
\begin{sheetcode}
from my_game import Grid  # no .py

def test_walls_block_moves():
    grid = Grid("###\n#P#\n###")
    assert grid.is_open(1, 1)
    assert not grid.is_open(0, 1)

def test_save(tmp_path):  # temp folder
    path = tmp_path / "scores.txt"
\end{sheetcode}
Run tests: Thonny \menu{Tools} $\rightarrow$ \menu{Open system shell\ldots}, go to your folder with
\texttt{cd}, then type \texttt{python3 -m pytest -q}.

\subsection{Before you commit: review checklist}
\begin{checklist}
  \item It runs, and the tests pass.
  \item Clear names; one job per function.
  \item Comments explain why.
  \item No real names, passwords or photos.
  \item Kind feedback: two stars and a wish.
\end{checklist}

\end{multicols}
\end{document}
